\documentclass{article}
\usepackage{amsmath}
\begin{document}


\section{Introduction}\label{sec:intro}
This is the first sentence.
This is the second sentence on the same line.
The penalty spring stiffness $k_p$ must be chosen carefully.
Display math should use double dollars: $$E = mc^2.$$
An inline equation like $F = ma$ appears here.

The method is described in~\cref{sec:method} and results in~\ref{sec:results}.
See also Table 1 and Figure 3 for details.


\section{Method}\label{sec:method}
A force of 100\,kN is applied at 34\,m/s.
The beam is 60\,mm long and weighs 500\,kg.

Use of \textbf{bold text} and \textit{italic text} is old-style.
Also \emph{emphasised} and \texttt{monospaced} should be updated.

\begin{equation}
F = ma,
\end{equation}
%
where $F$ is force.

\begin{equation}
E = \frac{1}{2}mv^2.
\end{equation}
%
The energy is always positive.


\section{Results}\label{sec:results}
The results show convergence for $n > 10$.
The range is from 0 to $\infty$\dots

\begin{figure}[t]
	\centering
	\includegraphics{fig1.pdf}
	\caption{A figure.}
	\label{fig:one}
\end{figure}

\begin{itemize}
	\item First item.
	\item Second item with $x = 1$.
\end{itemize}

We use \bigg( \frac{a}{b} \bigg) for large brackets.

\begin{enumerate}
	\item Alpha.
	\item Beta.
\end{enumerate}


\section{Conclusion}
The method works well.
Further work is needed. See Eq. 3 and Section 2 for details.
\end{document}
